\section{A common geometric dictionary, with integral models}
\label{sec:common-geometry}

Every restatement in this part retains the input encoding and the hypotheses
of Part~I.  A complexity theorem is not an isomorphism-invariant assertion
about an unencoded abstract scheme.  Published EGA/SGA/FGA locators support
the mathematical constructions.  The archive cotes identify inspected
historical companions; they are not presented as manuscripts containing AM,
ETH, or GCT complexity assertions.

\subsection{CX-G01: the representation-ring form of the analytic gates}
\label{finding:CX-G01}

\textbf{Finding.} Let $U=\operatorname{Spec}\mathbb Z[1/(B\ell)]$ and let
$\pi:Y\to U$ be the finite étale cover obtained by spreading out a number
field $E/\mathbb Q$, with all ramified primes inverted.  Its permutation
local system has class
\[
 [\pi_*\mathbb Q_\ell]
 = [\operatorname{Ind}_{G_E}^{G_{\mathbb Q}}\mathbf1]
 \quad\text{in }K_0(\operatorname{Rep}_{\mathbb Q_\ell}^{\mathrm{fin}}G_{\mathbb Q}).
\]
Its unramified Euler product is $\zeta_E$ with the excluded local factors
removed.  This exactly transports CX-R01--R02 and CX-K01--K04 into the
language of finite étale covers and representation Grothendieck groups.

\begin{proof}
The geometric fiber is the finite continuous $G_{\mathbb Q}$-set
$\operatorname{Hom}_{\mathbb Q}(E,\overline{\mathbb Q})$.
The stalk of the pushforward is the vector space of functions on that set.
A Frobenius cycle of length $f$ has characteristic polynomial $T^f-1$,
and hence reciprocal local factor $1-t^f$.  Its cycle lengths are the
residue degrees of the primes of $E$ above $p$.  Taking $t=p^{-s}$ gives
exactly the factors $(1-p^{-fs})$ in the reciprocal of the Dedekind Euler
factor.  Trace counts fixed points, hence degree-one primes, as needed by
Koiran.  The equivalence between finite étale covers and finite
fundamental-group sets is \cite[Exp.~V, \S4 Theorem~1, and \S7]{SGA1}.

For an abelian $E/\mathbb Q$, character decomposition gives the primitive
Dirichlet factorization in CX-R01.  For the cubic field of CX-R02, the
commutator trace calculation proves that the class is not a virtual sum of
rank-one objects, even allowing other conductors.  Taking $K_0$ does not
remove this obstruction.  The nonabelian Dedekind hypothesis remains the
same conditional input $M$.
\end{proof}

A finite-order ray-class character over $K$ instead defines a rank-one
local system on an open of $\operatorname{Spec}\mathcal O_K$, through
class field theory \cite[Introduction, Theorems~0.4--0.5]{MilneCFT}.
Its pushforward to an open of $\operatorname{Spec}\mathbb Z$ generally has
higher rank.  This is the exact dictionary for the Hecke hypotheses in
CX-N05--N08 and the class-field applications.  The corresponding objects
are Artin motives, with a coefficient field containing the finite character
values.  No general motivic conjecture is used.

\paragraph{Text and archive locations.}
Equivariant modules are treated in \cite[Exp.~I 4.7.1--4.7.2]{SGA3}.
The local determinant formalism is
\cite[Exp.~XV, \S3 no.~1, Proposition~1, and no.~2]{SGA5}; that source's
global rationality statement is over finite fields, not over
$\operatorname{Spec}\mathbb Z$.
The archival companions are cote~98, scan~12 (numbered sheet~11),
``Exemple I'' on finite $G$-sets, and scan~11 (sheet~10), the degree axioms
for norm morphisms \cite{Archive98}.

\subsection{CX-G02: multiplicative groups, Kummer torsors, and factoring}
\label{finding:CX-G02}

\textbf{Finding.} The modular operations of CX-N01--N06 are operations
on $\mathbb G_m$ and on explicit finite étale algebras.
An $r$-th power nonresidue $a\in\mathbb F_q^\times$, with $r\mid q-1$,
is precisely a nonzero class in
$\mathbb F_q^\times/(\mathbb F_q^\times)^r$, equivalently the Kummer
torsor
\[
 \operatorname{Spec}\mathbb F_q[t]/(t^r-a)
\]
has no $\mathbb F_q$-section.  A primitive root generates
$\mathbb G_m(\mathbb F_q)$; a generating set is a different output.

\begin{proof}
$\mathbb G_m=\operatorname{Spec}\mathbb Z[t,t^{-1}]$ represents units,
so its multiplication and power maps are literally the modular operations.
Miller's gcd tests take place in the same explicit residue ring
$\mathbb Z/N\mathbb Z$.
For the torsor, $r$ is invertible and $a\ne0$, so the derivative
$rt^{r-1}$ is a unit in the quotient.  The algebra is finite étale and
$\mu_r$ acts by multiplying $t$.  After adjoining a root, all other roots
differ from it by $\mu_r$, proving the torsor property.  A rational section
is exactly an $r$-th root of $a$.  Cyclicity gives the original power test
$a^{(q-1)/r}\ne1$ for prime $r$.  Primary-component assembly and the order
tests in Part~I are unchanged; the dictionary supplies no factorization of
the integer $q-1$.

For squarefree $f\in\mathbb F_q[t]$, the algebra
$A=\mathbb F_q[t]/(f)$ is finite étale.  Its product decomposition into
fields is the connected-component decomposition of $\operatorname{Spec}A$.
The Chinese remainder theorem sends a monic irreducible factor to its
primitive projector $e$.  Conversely, $(1-e)A$ has quotient the corresponding
field, and its inverse image in $\mathbb F_q[t]$ is generated by that factor.
Both conversions use polynomial arithmetic and linear algebra.  Repeated
factors are treated by squarefree decomposition, including Frobenius-root
extraction over the perfect base field.  Thus factoring is the effective
idempotent decomposition problem.  Existence of the finite étale
decomposition does not compute its idempotents or change the bound of the
selected factoring algorithm.
\end{proof}

\paragraph{Locations.}
$\mathbb G_m$ and $\mu_r$ occur in the diagonalizable-group formalism
\cite[Exp.~VIII 1.1--1.2]{SGA3}; the exact torsor sequence is
\cite[Exp.~IX 3.2]{SGA4}.  See also \cite[Exp.~V \S7]{SGA1}.
Cote~28, scan~3 (sheet~2), writes the action map $G\times X\to X$ and its
axioms \cite{Archive28}; cote~98, scan~12 (sheet~11), supplies the finite
$G$-set companion \cite{Archive98}.  These are locations of the language,
not historic primality or factoring algorithms.

\paragraph{The extension refinements in this language.}
CX-N13 constructs a nontrivial class in
$H^1_{\mathrm{et}}(\operatorname{Spec}\mathbb F_{p^m},\mu_r)$
for a specified prime $r\mid p-1$ from the supplied field model.
Its odd-$r$ Frobenius eigenspace is computed on that finite free
$\mathbb F_p$-algebra; its power test is exactly the torsor's
failure to have a section. The even-degree quadratic case retains
its unconditional construction.
For CX-N14, a supplied automorphism of a constant finite étale
$\mathbb F_p$-scheme gives the same component action. The moved-component
ideal and the descended idempotent construct a proper open-and-closed
subscheme. The cases $p\le d$, $r\mid p-1$, and
$v_r(p^{\operatorname{ord}_r(p)}-1)=1$, and the optional coefficient
Teichmüller unit, retain precisely their hypotheses and bounds.
For CX-N15, the Frobenius-fixed algebra
$B_f=\ker(a\mapsto a^p-a)$ inside
$A=\mathbb F_q[X]/(f)$ is $\mathbb F_p^t$, where $t$ is the number
of connected components of the finite étale scheme $\operatorname{Spec}A$.
Its primitive idempotents are exactly those of $A$.
The coefficient extension, eigenspace and equivariant interpolation
construct the unit needed to split this constant component scheme.
The resulting open-and-closed decompositions of $\operatorname{Spec}A$
therefore compute the factors in
$(d^{\log d}\log q)^{O(1)}$ time, under the exact cyclotomic--Kummer
Dedekind-zeta family in CX-N15. The pointwise Frobenius calculation is
an algebraic construction; its required number-field prime bound is
still the stated conditional analytic input.
Thus the Kummer and descent operations carry the complete proofs,
not merely their existence conclusions, into this dictionary.

\subsection{CX-G03: ideal-class computation as Picard-group computation}
\label{finding:CX-G03}

\textbf{Finding.} For $X=\operatorname{Spec}\mathcal O_K$,
the class group is $\operatorname{Pic}(X)$.  CX-N07 bounds the prime-divisor
line bundles generating this group, with exactly the same conditional
premise.  CX-N08 and the concrete class-group computation concern completeness
of the principal-divisor and unit data, with their original search bounds.

\begin{proof}
A nonzero fractional ideal $I$ of the Dedekind domain $\mathcal O_K$ is an
invertible module.  Tensor product corresponds to ideal multiplication.
Conversely, choose a basis of the generic fiber of a rank-one locally free
module and embed it in $K$; clearing denominators identifies it with a
fractional ideal.  An isomorphism of two such modules extends generically
to multiplication by a scalar in $K^\times$.  Thus isomorphism classes
are exactly fractional ideals modulo principal fractional ideals.
The prime ideal $\mathfrak p$ corresponds, in effective-divisor notation,
to $\mathcal O_X(-\mathfrak p)$; changing all signs does not affect generation.
The factor-base map is the divisor-to-line-bundle map, with kernel the
principal-divisor relations supported on that base.  Units are the kernel
of the principal-divisor map from $K^\times$.  These are precisely the
objects in the exact sequence behind CX-N08, so its index identity and
its separate relation-search cost are preserved.
\end{proof}

The Hilbert-class-field cover is
$\operatorname{Spec}\mathcal O_{H_K}\to\operatorname{Spec}\mathcal O_K$,
finite étale at finite places.  For ray class fields, delete the modulus
primes and retain the real-place conditions in the arithmetic character data.
Their characters give the local systems of CX-G01.  The required analytic
factor is $\zeta_{H_K}$ or the pertinent Hecke factors, not automatically
$\zeta_K$ alone.

\paragraph{Locations.}
The cohomological identification is
\cite[Exp.~IX, Theorem~3.3]{SGA4}; its $K_0$ and $\gamma$ language is in
\cite[Exp.~V 3.10; Exp.~VI 3.2--3.3]{SGA6}.
Cote~2, scan~5 (sheet~4), discusses rank augmentation and Picard comparison
in dimension at most one \cite{Archive2}.  No class-group complexity
theorem is attributed to that manuscript.

\subsection{CX-G04: integral orbit closures and the role of Hilbert and Quot schemes}
\label{finding:CX-G04}

\textbf{Finding.} Let $G=\operatorname{GL}_{m^2,\mathbb Z}$ act on the
coefficient space $V_{\mathbb Z}$ of degree-$m$ forms in $m^2$ variables.
If $I_{\mathbb Q}$ is the ideal of the characteristic-zero determinant
orbit closure, then
\[
 I_{\mathbb Z}=I_{\mathbb Q}\cap\mathbb Z[V],\qquad
 X_{\mathbb Z}=\operatorname{Spec}(\mathbb Z[V]/I_{\mathbb Z})
\]
defines a flat $G$-scheme over $\mathbb Z$.  Its projectivization is a flat
projective family and defines a point of the appropriate Hilbert scheme.
This supplies integral objects for CX-E06--E08 and CX-K05--K08, without
asserting special-fiber equality with every characteristic-$p$ orbit closure.

\begin{proof}
The quotient embeds in $\mathbb Q[V]/I_{\mathbb Q}$ and is torsion free
over $\mathbb Z$.  A torsion-free module over the PID $\mathbb Z$ is a
filtered union of finitely generated torsion-free submodules, which are
free.  Filtered colimits preserve flatness, proving flatness.

The determinant orbit closure is a cone in characteristic zero: over an
algebraic closure, any nonzero scalar multiple of the determinant is
obtained by scalar change of variables.  Its ideal is therefore homogeneous,
as is its contraction.  Every graded quotient piece is finitely generated
and torsion free over $\mathbb Z$, hence free.  On a standard affine chart
of Proj, localization preserves flatness and the degree-zero component is
a direct summand over $\mathbb Z$.  Thus Proj is flat.  It is finitely
presented by noetherianity.  The graded pieces have the same dimensions in
every fiber, so the fiber Hilbert polynomials coincide.  For their polynomial
$P$, \cite[\S3, Theorem~3.2]{FGA221} supplies the classifying morphism
\[
 \operatorname{Spec}\mathbb Z\longrightarrow
 \operatorname{Hilb}^{P}(\mathbb P(V)_{\mathbb Z}/\mathbb Z).
\]

The integral change-of-variables formula defines the original coaction.
For $f\in I_{\mathbb Z}$ its image in
$\mathbb Z[G]\otimes(\mathbb Z[V]/I_{\mathbb Z})$ vanishes after tensoring
with $\mathbb Q$.  This target is flat and torsion free over $\mathbb Z$,
so the image already vanishes.  Hence the ideal is $G$-stable, proving the
asserted integral action.  This argument does not claim that formation of
a scheme-theoretic image commutes with all special-fiber base changes.
\end{proof}

Neither finite presentation nor representability gives an efficient list
of equations or a succinct computation of $P$.  No step bounds the degrees
or bit sizes of ideal generators as a function of the succinct circuit
parameter.  The independent equation-access gate of CX-K06 therefore remains.
The effective statement is the search for a finite morphism in CX-K08;
placing its data in Hilbert or Quot schemes does not compute that search.

\paragraph{Locations.}
Linear spaces, actions, and comodules are
\cite[Exp.~I 4.6.3.1, 4.7.1--4.7.2]{SGA3}; finite morphisms and base change
are \cite[6.1.4--6.1.5]{EGA2}.  The parameter spaces and their projective
embeddings are \cite[\S3, Theorems~3.1--3.2, pp.~259--260;
Proposition~3.8, p.~265]{FGA221}.
The archival companions are cote~28, scan~3 (sheet~2), and cote~98,
scan~44 (sheet~43), headed ``Familles limit\'ees de faisceaux coh\'erents''
\cite{Archive28,Archive98}.  The latter sheet is not used in place of the
published representability theorem.

\subsection{CX-G05: exterior powers, Schur operations, and loss of information under rank}
\label{finding:CX-G05}

\textbf{Finding.} The representation objects and locally free sheaves above
have exact-category Grothendieck rings, tensor multiplication, and exterior
operations $\lambda^i$.  With rank augmentation $\varepsilon$, the usual
operations on rank-zero classes are
\[
 \gamma_t(x)=\lambda_{t/(1-t)}(x),\qquad
 F^j_\gamma=\left\langle
  \prod_a\gamma^{i_a}(x_a):\ \varepsilon(x_a)=0,\ \sum_a i_a\geq j
 \right\rangle .
\]
These are the existing structures of
\cite[Exp.~V 2.4, 3.10(3.10.1); Exp.~VI 3.2--3.3]{SGA6}.
For the Boolean scheme of CX-E07, with $N=\#F$,
\[
 [p_*\mathcal O_{X_F}]=N[\mathcal O],\qquad
 \lambda_t(N[\mathcal O])=(1+t)^N.
\]
Its reduced rank class is zero.  This rank data does not determine the
equivariant coordinate-ring classes required for the obstruction in CX-E08.

\begin{proof}
CX-E07 identifies the pushforward with a direct sum of $N$ copies of
$\mathcal O$.  Its $i$-th exterior power has rank $\binom Ni$, giving the
displayed series.  Subtracting $N[\mathcal O]$ gives zero, so all positive
$\gamma$-operations on that reduced class vanish, functorially after base
change.

For the loss of representation data, the trivial and sign representations
of $C_2$ over $\mathbb Q$ both have rank one.  Their traces at the nonidentity
element are $1$ and $-1$.  Additivity of trace in exact sequences shows that
their $K_0$ classes differ.  Thus rank is not injective.  A count transported
only by the rank-class dictionary cannot specify the equivariant class,
let alone its graded Schur multiplicities.  This refutes that particular
reconstruction from rank; it is not a complexity separation or a claim
that no extra reduction could ever provide such data.
\end{proof}

Over a $\mathbb Q$-scheme, Schur functors on locally free $\mathcal E$ are
images of rational Young idempotents on $\mathcal E^{\otimes d}$.
Idempotents split, giving locally free images.  One can likewise work after
inverting the required factorial denominators.  This is not an assertion
of integral semisimplicity over $\mathbb Z$.
Their classes are universal Schur operations; integrally these classes can
be written using Jacobi--Trudi,
\[
 s_\lambda(\mathcal E)
 =\det(h_{\lambda_i-i+j}(\mathcal E)),\qquad
 \sum_{i\geq0}h_i(\mathcal E)t^i=\lambda_{-t}([\mathcal E])^{-1},
\]
with $h_0=1$ and $h_i=0$ for $i<0$
\cite[Chapter~I \S3, (3.4), p.~41]{Macdonald1995}.
The representation interpretation is the equivariant-module language of
\cite[Exp.~I 4.7.1--4.7.2]{SGA3}.  Cote~2, scan~5 (sheet~4), is the rank
augmentation companion, and cote~98, scan~12 (sheet~11), is the finite
$G$-set companion \cite{Archive2,Archive98}.  No plethysm or Kronecker
problem from the parallel prompts is being asserted solved.

\subsection{CX-G06: the exact boundary between pointwise weights and global GRH}
\label{finding:CX-G06}

\textbf{Finding.} The finite-monodromy arithmetic sheaves in CX-G01 have
root-of-unity Frobenius eigenvalues at good primes, hence pointwise weight
zero.  Tensor, exterior, and characteristic-zero direct-summand operations
preserve that property.  This fact does not supply an additional proof of
the global zero-location input used in Part~I.

\begin{proof}
A finite-order matrix is annihilated by $T^a-1$ for some $a$; all its
eigenvalues are roots of unity.  Tensor eigenvalues and exterior-power
eigenvalues are products of these, and eigenvalues of a stable direct
summand occur among those of the original representation.  All complex
absolute values are one.  This proves precisely the pointwise assertion.

The number-field GRH assertion concerns analytic continuation and zeros of
an infinite Euler product in $s$.  The preceding finite-order linear-algebra
argument supplies no such continuation or zero estimate.  In our proved
implications the required global assertion remains $A_D$ or $M$, exactly
as recorded in CX-G01.  This identifies the scope of the argument, not
logical independence of mathematical theories.

For the Boolean scheme, CX-E07 gives the stronger explicit description.
After base change to $\overline{\mathbb F}_q$, its étale cohomology is $N$
copies of $\mathbb Q_\ell$ in degree zero and zero in higher degrees.
Frobenius is the identity.  Hence the fiber zeta function is $(1-t)^{-N}$
and its Artin motive is a sum of $N$ unit motives.  The counting algorithm
computes this multiplicity from a succinct presentation; it constructs
no nontrivial motive and changes no purity theorem.
\end{proof}

The existing finite-field results are
\cite[\S1.2; Theorem~3.3.1, p.~204; Corollary~3.3.9, p.~207]{WeilII},
with the Euler-factor and direct-image formalism of
\cite[Exp.~XV \S3, no.~1 Proposition~1, and no.~2]{SGA5}.
The inspected archival comparison is cote~34, scan~64 (sheet~63),
crossed-out Remark~1.17, which relates a proposed openness assertion to
Riemann--Weil in higher dimension; scan~74 (sheet~73), equation~(2.5.5),
uses the curve case for an invariant-cohomology vanishing
\cite{Archive34}.  The crossed-out status of the first passage is part
of the locator.  These are finite-field cohomological companions, not
number-field zero-location statements.  Our finding changes neither
their mathematical scope nor the published Weil~II theorem.

\section{Where every finding lives in this dictionary}
\label{sec:crosswalk}

\begin{longtable}{@{}p{0.17\textwidth}p{0.42\textwidth}p{0.34\textwidth}@{}}
\toprule
Findings & Restatement and proof-carrying dictionary & Exact locations and archival companions\\
\midrule
\endhead
CX-R01--R02 & The finite étale permutation object, rank-one decomposition, and commutator obstruction of CX-G01. & SGA1 V \S4 Thm.~1, \S7; SGA3 I4.7.2; cote98/12 (sheet11).\\
CX-K01--K04 & Nonempty generic fiber, bounded closed point, degree-one specializations; abelian promise and cubic obstruction. CX-K08 and CX-G01 carry the proofs. & SGA1 V \S7; SGA5 XV \S3 no.1; cote98/11--12 (sheets10--11).\\
CX-K05--K07 & Search for a finite map from a graded affine scheme to affine space; a small separating system and a dimension-size system retain different output requirements. CX-K08 and CX-G04 preserve these distinctions. & EGAII6.1.4--6.1.5; FGA221 Thms3.1--3.2; cote98/44 (sheet43).\\
CX-K08 & Integral spreading, finite étale cover, Frobenius trace, and open locus of finite linear projections. & SGA1 V \S7; EGAII6.1.4; FGA221 \S3; cote28/3 and cote98/11--12,44.\\
CX-N01--N04 & Explicit residue-ring and $\mathbb G_m$ computations; Kummer torsor classes and generators, retaining the factorization input. CX-G02. & SGA3 VIII1.1--1.2; SGA4 IX3.2; cote28/3, cote98/12.\\
CX-N05--N06 & Ray-character local systems on arithmetic curves and their direct images, together with finite-field group computations. CX-G01--G02. & SGA1 V \S7; SGA3 I4.7.2; SGA4 IX3.2; cote98/11--12.\\
CX-N07--N08, CX-N12 & Picard generators, $S$-units, principal-divisor relations, and the same certification index. CX-G03. & SGA4 IX3.3; SGA6 V3.10, VI3.2; cote2/5 (sheet4).\\
CX-N09--N11, CX-N13--N15 & Connected-component/idempotent computation for finite étale algebras, Kummer power classes, and Frobenius in the stated arithmetic cover. CX-G01--G02; the prime-field automorphism and general composition proofs construct open-and-closed decompositions with their original bounds. & SGA1 V \S7; SGA4 IX3.2; SGA3 VIII1.1; cote98/12, cote28/3.\\
CX-E01--E05 & Emptiness and degree of the succinct Boolean finite étale scheme, parsimonious gate isomorphisms, and the original Boolean circuit model. The table and proof in CX-E07 give each restatement. & SGA1 V \S7; SGA6 VI3.2--3.3; SGA5 XV \S3; cote98/11--12, cote2/5.\\
CX-E06, CX-E08 & Symbolic polynomials and graded equivariant coordinate rings of the integral orbit models; multiplicity quotient condition. CX-G04--G05. & SGA3 I4.7.2; FGA221 Thm3.2; SGA6 V2.4; cote28/3, cote98/44.\\
CX-E07 & Degree equals rank and the multiplicity of the unit Artin motive; reduced $\gamma$-class is zero. CX-G05--G06. & SGA6 V3.10, VI3.2; SGA5 XV \S3; cote2/5, cote98/11.\\
CX-G01--G06 & The displayed geometric proofs are already in the requested language; no further translation changes their premises. & The individual exact locators and scan numbers are supplied with each finding.\\
\bottomrule
\end{longtable}

Here ``cote98/12'' means PDF scan~12 of cote~98, including the archival
cover as scan~1; the sheet itself is numbered~11.  The same convention is
used throughout.  The archive references are to the inspected page content,
not a claim of dating or priority beyond the archive's catalogue.
